\section{Определение матрицы перехода от базиса к базису}

\begin{definition}[Матрица перехода от базиса к базису]
Пусть $V$ --- линейное пространство над полем $F$.

$E_1 = \{e_1, \ldots, e_n\}$ --- старый базис.

$E_2 = \{e'_1, \ldots, e'_n\}$ --- новый базис.

Разложим новые базисные векторы по старому базису:
\begin{align*}
e'_1 &= t_{11}e_1 + \cdots + t_{n1}e_n, \\
&\vdots \\
e'_n &= t_{1n}e_1 + \cdots + t_{nn}e_n.
\end{align*}

Для любого $x \in V$: $x = \xi E_1$ и $x = \xi' E_2$, где $\xi, \xi'$ --- координаты вектора.

Матрица
\[
T = \begin{pmatrix}
t_{11} & \cdots & t_{1n} \\
\vdots & \ddots & \vdots \\
t_{n1} & \cdots & t_{nn}
\end{pmatrix}
\]
называется \textbf{матрицей перехода} от базиса $E_1$ к базису $E_2$.

В матричной форме: $E_2 = E_1 T$.
\end{definition}
